% Core
\newglossaryentry{imaginary_constant}{
    name={imaginary constant},
    description={The square root of -1, written $i$},
    plural={imaginary constants}
}
\newglossaryentry{real_number}{
    name={real number},
    description={A number on the real number line; a number in the interval $(-\infty, \infty)$.
    The set of real numbers is written $\R$},
    plural={real numbers}
}
\newglossaryentry{complex_number}{
    name={complex number},
    description={A number in the form $a+bi$, where $a,b\in\R$.
    The set of complex numbers is written $\C$},
    plural={complex numbers}
}
\newglossaryentry{polynomial}{
    name={polynomial},
    description={An expression in the form $a_0+a_1x^1+a_2x^2+\cdots+a_nx^n$},
    plural={polynomials}
}
\newglossaryentry{complex_conjugate}{
    name={complex conjugate},
    description={A complex number with a flipped sign before its imaginary part.
    The complex conjugate of a complex number $z$ is written $z^*$},
    plural={}
}
\newglossaryentry{root}{
    name={root},
    description={A value that satisfies an equation, for example 1 and -1 are the roots of the equation $x^2-1=0$, since for $x=1$ and $x=-1$, the equation holds},
    plural={roots}
}
\newglossaryentry{complex_root}{
    name={complex root},
    description={A root $x$ of an equation where $x\in\C$.
    It is often also used to mean that $x\notin\R$},
    plural={complex roots}
}
\newglossaryentry{cartesian_graph}{
    name={Cartesian graph},
    description={A graph in the $x$-$y$ plane},
    plural={Cartesian graphs}
}
\newglossaryentry{realising_the_denominator}{
    name={realising the denominator},
    description={The process of multiplying a complex fraction in the form $\frac{w}{z}$ by $\frac{z^*}{z^*}$ (where $w,z\in\C$) to make the denominator a real value without changing the value of the fraction.},
}
\newglossaryentry{quadratic}{
    name={quadratic},
    description={A second degree \gls{polynomial}, i.e. an equation in the form $ax^2+bx+c=0$, where $a\neq0$.
    When plotted on a \gls{cartesian_graph}, it forms a parabola.
    Its \glspl{root} are given by the quadratic formula, $x=\frac{-b\pm\sqrt{b^2-4ac}}{2a}$},
    plural={quadratics}
}
\newglossaryentry{conjugate_pair}{
    name={conjugate pair},
    description={A pair of \glspl{complex_number}, where one is the \gls{complex_conjugate} of the other},
    plural={conjugate pairs}
}
\newglossaryentry{cubic}{
    name={cubic},
    description={A 3rd degree \gls{polynomial}, i.e. an equation in the form $x^3+bx^2+cx+d=0$, where $a\neq0$},
    plural={conjugate pairs}
}
\newglossaryentry{argand_diagram}{
    name=Argand diagram,
    description={A 2D graph of the complex plane},
    plural={Argand diagrams}
}
\newglossaryentry{real_part}{
    name={real part},
    description={The real term of a \gls{complex_number}, i.e. where $z=a+bi,\ a,b\in\Relbar$, the real part of $z$ is $a$
    The real part of a \gls{complex_number} $z$ is written as $\real(z)$ or $\Re(z)$},
    plural={}
}
\newglossaryentry{imaginary_part}{
    name={imaginary part},
    description={The coefficient of the imaginary term of a complex number, i.e. where $z=a+bi,\ a,b\in\R$, the imaginary part of $z$ is $b$.
    The imaginary part of a \gls{complex_number} $z$ is written as $imag(z)$ or $\Im(z)$},
    plural={}
}
\newglossaryentry{modulus_argument_form}{
    name={modulus-arguemnt form},
    description={The form $r(\cos\theta+i\sin\theta)\equiv\cis\theta$},
    plural={}
}
\newglossaryentry{radian}{
    name={radian},
    description={An angle unit, equal to the angle required to make the arc length of a unit circle equal 1.
    A whole circle is $2\pi$ radians; half a circle is $\pi$ radians.
    The symbol for radians is either $^\text{c}$ or $\text{rad}$, as in $\pi^\text{c}$ or $\pi\ \text{rad}$, or is often omitted},
    plural={radians}
}
\newglossaryentry{modulus}{
    name={modulus},
    description={The magnitude of a number, where $|x|$ is the modulus of $x$.
    For positive \glspl{real_number}, the modulus is itself, for negative \glspl{real_number}, the modulus is the positive version of itself ($|x|=-x,\ x\in\R^-$).
    For a \gls{complex_number} $z$, $|z|=\sqrt{\real^2(z)+\imag^2(z)}$}, % TODO: come back
    plural={}
}
\newglossaryentry{quadrant}{
    name={quadrant},
    description={A quarter of a graph; the 1st quadrant is the top right, the second is the top left, the third is the bottom left, and the fourth is the bottom right.},
    plural={quadrants}
}
\newglossaryentry{argument}{
    name={argument},
    description={The angle that a \gls{complex_number} makes with the positive real number line.
    It is written as $\arg(z)$.
    The argument is in the real interval $(-\pi,\pi]$},
    plural={arguments}
}
\newglossaryentry{integer}{
    name={integer},
    description={A number with no decimal part.
    The set of all integers \\ ($\{\cdots-3,-2,-1,0,1,2,3,\cdots\})$ is written as $\Z$},
    plural={arguments}
}



% Further Mechanics
\newglossaryentry{mass}{
    name={momentum},
    description={},
    plural={masses},
    symbol={$m$}
}
\newglossaryentry{velocity}{
    name={momentum},
    description={Velocity is a vector quantity measuring the rate of change of position of an object.
    See equations \ref{}}, % TODO: Put eqn refs here for definitions of velocity
    plural={velocities},
    symbol={\textbf{v}}
}
\newglossaryentry{momentum}{
    name={momentum},
    description={The product of an object's \gls{mass} and \gls{velocity}.
    It the \glspl{si_unit} \gls{kilogram} \glspl{metre} per \gls{second} (kgms$^{-1}$), or \gls{newton} \glspl{second} (Ns).
    See equations \ref{eq:momentum_m_v}, \ref{eq:momentum_f_t}},
    plural={momenta},
    symbol={\textbf{p}}
}
\newglossaryentry{force}{
    name={force},
    description={A push or pull on an object.
    It has the \glspl{si_unit} \glspl{newton} (N).
    See equations \ref{}}, % TODO: Put eqn refs here for definitions of force
    plural={forces},
    symbol={\textbf{F}}
}
\newglossaryentry{time}{
    name={time},
    description={The progression of existence, represented as the fourth dimension of (3+1)D spacetime.
    It has the \glspl{si_unit} \glspl{second} (s)},
    plural={times},
    symbol={$t$}
}
\newglossaryentry{si_unit}{
    name={SI unit},
    description={The International System of units; the standard unit system used globally for maths and other STEM subjects.
    All SI units are made from the 7 base units: the \gls{metre} (m), \gls{kilogram} (kg), the \gls{second} (s), the \gls{ampere} (A), the \gls{kelvin} (K), the \gls{mole} (mol) and the \gls{candela} (cd).
    All of the base units are derived from 7 universal constants: the speed of light ($c$), the hyperfine transition frequency of caesium-133, the Planck constant ($h$), the Avrogadro constant ($N_A$), the luminous efficacy of 540THz radiation, the Boltzmann constant ($k$), and the elementary change ($e$)},
    plural={SI units}
}
\newglossaryentry{metre}{
    name={metre},
    description={The \gls{si_unit} of length},
    plural={metres},
    symbol={m}
}
\newglossaryentry{kilogram}{
    name={kilogram},
    description={The \gls{si_unit} of mass.
    It appears to have an SI prefix (kilo), however due to historical reasons, it is considered a base unit.
    The gram still exists, defined as a thousandth of a kilogram},
    plural={kilograms},
    symbol={kg}
}
\newglossaryentry{second}{
    name={second},
    description={The \gls{si_unit} of time},
    plural={seconds},
    symbol={s}
}
\newglossaryentry{ampere}{
    name={ampere},
    description={The \gls{si_unit} of electric current},
    plural={amperes},
    symbol={A}
}
\newglossaryentry{kelvin}{
    name={kelvin},
    description={The \gls{si_unit} of temperature.
    A change of temperature of 1K is the same as a change of 1$^\circ$C, but 0K is approximately -273$^\circ$C},
    plural={kelvin},
    symbol={K}
}
\newglossaryentry{mole}{
    name={mole},
    description={The \gls{si_unit} of quantity.
    1mol is approximately the number of atoms of carbon-12 in 12g of carbon-12.
    It is defined as $N_A^-1$},
    plural={moles},
    symbol={mol}
}
\newglossaryentry{candela}{
    name={candela},
    description={The \gls{si_unit} of luminous intensity.
    It is approximately the luminous intensity of a wax candle},
    plural={candela},
    symbol={cd}
}
\newglossaryentry{newton}{
    name={newton},
    description={The \gls{si_unit} of force.
    The newton is not a base unit, but instead is equivalent to a \gls{kilogram} \gls{metre} per \gls{second} squared (kgms$^{-2}$)},
    plural={newtons},
    symbol={N}
}



% Decision
\newglossaryentry{algorithm}{
    name={algorithm},
    description={A set of instructions instructing a person or computer on what steps to follow to achieve a goal.
    See section \ref{algorithms}},
    plural={algorithms}
}
\newglossaryentry{trace_table}{
    name={trace table},
    description={A table that tracks variables while an \gls{algorithm} is running},
    plural={trace tables}
}
\newglossaryentry{instruction}{
    name={instruction},
    description={A step of an \gls{algorithm}, telling the person/computer running the \gls{algorithm} what to do at a specific point},
    plural={instructions}
}
\newglossaryentry{sorting_algorithm}{
    name={sorting algorithm},
    description={A type of \gls{algorithm} that takes a list of items and sorts them.
    One example is \gls{bubble_sort}},
    plural={sorting algorithms}
}
\newglossaryentry{bubble_sort}{
    name={bubble sort},
    description={A \gls{sorting_algorithm} that looks through each pair of items and swaps them if they are in the wrong order.
    It has an order of $\bigO(n^2)$.
    See section \ref{bubble_sort}},
    plural={bubble sorts}
}
\newglossaryentry{graph}{
    name={graph},
    description={A set of \glspl{vertex}, connected by \glspl{edge}, also called a network.
    See section \ref{graphs_and_networks}},
    plural={graphs}
}
\newglossaryentry{vertex}{
    name={vertex},
    description={A point of a graph, also called a node},
    plural={vertices}
}
\newglossaryentry{edge}{
    name={edge},
    description={A line connecting two \glspl{vertex}, also called an arc.
    Edges can have weights (weighted edges) and directions (directed edges).
    These are used to make weighted graphs and directed graphs respectively},
    plural={edges}
}
\newglossaryentry{weighted_edge}{
    name={weighted edge},
    description={An \gls{edge} with a weight assigned to it.
    This is a scalar usually representing a real-world value.
    For example, if each \gls{edge} represents a road connecting two places (two \glspl{vertex}), then the weight could represent the length of the road, or the time taken to travel it},
    plural={weighted edges}
}
\newglossaryentry{directed_edge}{
    name={directed edge},
    description={An \gls{edge} that only connects its two \glspl{vertex} in one direction, and cannot be travelled in the other direction},
    plural={directed edges}
}
\newglossaryentry{weighted_graph}{
    name={weighted graph},
    description={A \gls{graph} containing one or more \glspl{weighted_edge}},
    plural={weighted graphs}
}
\newglossaryentry{directed_graph}{
    name={directed graph},
    description={A \gls{graph} containing one or more \glspl{directed_edge}},
    plural={directed graphs}
}
\newglossaryentry{vertex_degree}{
    name={degree},
    description={(Graph theory) The the number of \glspl{edge} connecting to a given \gls{vertex}},
    plural={directed graphs}
}
\newglossaryentry{walk}{
    name={walk},
    description={A route along \glspl{edge} of a \gls{graph}},
    plural={walks}
}
\newglossaryentry{path}{
    name={path},
    description={A \gls{walk} where no \gls{vertex} is visited twice},
    plural={paths}
}
\newglossaryentry{trail}{
    name={trail},
    description={A \gls{walk} where no \gls{edge} is visited twice},
    plural={trails}
}
\newglossaryentry{cycle}{
    name={cycle},
    description={A \gls{walk} where the first \gls{vertex} visited is also the last \gls{vertex} visited, and no other \gls{vertex} is visited twice},
    plural={cycles}
}
\newglossaryentry{hamiltonian_cycle}{
    name={Hamiltonian cycle},
    description={A \gls{cycle} where every \gls{vertex} the \gls{graph} is visited},
    plural={Hamiltonian cycles}
}



\newglossaryentry{UNUSED_WITHOUT_SYMBOL}{
    name={},
    description={},
    plural={}
}

\newglossaryentry{UNUSED_WITH_SYMBOL}{
    name={},
    description={},
    plural={},
    symbol={}
}